\chapter{Архитектурный анализ и проектирование системы анализа}
\label{ch:design}

\section{Теоретические предпосылки и требования к системе анализа}

Разработка прикладного инструмента для анализа вычислительных возможностей вычислительной платформы требует исходить из того, что производительность не является простой и однозначной характеристикой. На уровне пользовательского восприятия компьютер может рассматриваться как единое устройство, однако на уровне инженерного анализа его поведение определяется взаимодействием нескольких неоднородных подсистем. Итог любой вычислительной нагрузки зависит от эффективности выполнения инструкций в ядре процессора, степени использования параллелизма, скорости и латентности доступа к памяти, состояния дисковой и сетевой подсистем, а также от фонового воздействия операционной системы и активных процессов. Следовательно, проектирование системы анализа не может ограничиваться регистрацией одной-двух агрегированных метрик; необходимо формировать архитектуру, ориентированную на комплексный сбор, интерпретацию и сопоставление данных разных типов.

Теоретической предпосылкой проектирования является признание того факта, что для разных классов приложений «узкое место» формируется по-разному. Для одних задач критична вычислительная плотность и способность системы поддерживать высокий темп арифметических операций. Для других определяющим фактором становится память: либо из-за ограниченной пропускной способности при потоковой обработке данных, либо из-за высокой латентности случайного доступа. Для третьих нагрузок основным ограничением оказывается подсистема ввода-вывода, если приложение активно взаимодействует с накопителем, сетью или внешними устройствами. По этой причине архитектура инструмента анализа должна обеспечивать возможность наблюдать не только итоговый результат теста, но и тот контекст, в котором этот результат был получен.

Для вычислительных нагрузок критичен параллелизм как на уровне инструкций, так и на уровне потоков исполнения. Современные процессоры используют внеочередное исполнение, конвейеризацию, предсказание ветвлений, многоуровневые кэши и, во многих случаях, гетерогенную организацию ядер. В таких условиях одна и та же программа может демонстрировать существенно различающуюся эффективность в зависимости от структуры кода, числа потоков, локальности доступа к данным и поведения планировщика. Именно поэтому в проектируемом комплексе предусмотрены тесты, формирующие различный профиль процессорной нагрузки: сортировка большого массива данных нагружает кэш-подсистему и предсказатель ветвлений; рекурсивные вычисления акцентируют накладные расходы вызовов функций, работу со стеком и чувствительность к последовательной логике; параллельные вычисления по методу Монте-Карло позволяют оценивать масштабируемость по числу доступных ядер и эффективность конкурентного исполнения.

Подсистема памяти выступает самостоятельным источником ограничений, который часто проявляется менее очевидно, чем недостаток вычислительной мощности. Во многих приложениях процессор простаивает не потому, что не способен выполнять операции достаточно быстро, а потому, что ожидает поступления данных из кэша более низкого уровня или из основной памяти. По этой причине в проекте реализованы два принципиально разных класса измерений памяти: оценка пропускной способности линейных операций и оценка латентности нерегулярных обращений. Первый тип измерений отражает возможности контроллера памяти и общую эффективность работы с последовательными массивами данных. Второй тип ориентирован на имитацию практических сценариев со слабой пространственной локальностью, где используется случайная цепочка индексов по принципу pointer-chasing, снижающая влияние аппаратной предвыборки и позволяющая приблизиться к более реалистичной оценке задержки доступа.

В контексте компьютерной безопасности существенную роль играет и криптографическая производительность. Современные системы постоянно выполняют шифрование, хеширование, аутентификацию, проверку цифровых подписей и генерацию ключевого материала. Следовательно, оценка вычислительной платформы без учёта криптографических операций была бы неполной, особенно если предполагается использование системы в сетевых сервисах, системах защищённого хранения данных или инфраструктуре аутентификации. В проекте поэтому выделен отдельный набор криптографических измерений. С одной стороны, алгоритмы хеширования, такие как SHA-256 и SHA-512, характеризуют скорость проверки целостности и обработки потоков данных. С другой стороны, симметричные и асимметричные алгоритмы, включая AES, RSA, Ed25519 и ChaCha20-Poly1305, позволяют судить о пропускной способности защищённых каналов и стоимости операций, характерных для протоколов безопасности. Для этой категории в системе реализованы как автоматические тесты, так и интерактивный режим, позволяющий изменять параметры и наблюдать, как они влияют на скорость.

Таким образом, уже на этапе формулирования требований становится очевидным, что разрабатываемая система анализа должна опираться на комбинированный подход. С одной стороны, она должна собирать и отображать телеметрию в реальном времени, фиксируя состояние платформы «здесь и сейчас». С другой стороны, она должна обеспечивать воспроизводимые измерения с использованием синтетических тестов, нацеленных на разные классы нагрузок. Только при объединении этих двух составляющих пользователь получает возможность не просто наблюдать показатели, а интерпретировать их с учётом причин, условий и ограничений.

\subsection{Функциональные требования}

С точки зрения пользовательских сценариев комплекс должен обеспечивать несколько взаимосвязанных групп функций. Первая группа связана с наблюдением за текущим состоянием системы. Пользователь должен получать информацию о загрузке процессора, распределении активности по логическим ядрам, состоянии памяти, дисковой и сетевой активности, времени работы системы, средней нагрузке и перечне процессов, формирующих наибольшее потребление ресурсов. При этом отображение должно быть организовано таким образом, чтобы метрики воспринимались не как набор изолированных чисел, а как связанная картина состояния платформы.

Вторая группа функций относится к синтетическому тестированию. Система должна позволять запускать измерительные сценарии по нескольким категориям нагрузок, соответствующим ключевым подсистемам вычислительной платформы. Пользователь должен иметь возможность инициировать тест, получить итоговый результат в естественных прикладных единицах измерения, ознакомиться с пояснением смысла полученного числа и сопоставить его с качеством измерения. Последний аспект является особенно важным, поскольку без информации о разбросе, выбросах или коэффициенте вариации числовой результат остаётся потенциально неоднозначным.

Третья группа функций связана с интерпретацией и сопровождением результата. Пользовательский интерфейс должен не только выводить итоговое значение, но и пояснять, как трактовать его с прикладной точки зрения. Например, высокая пропускная способность памяти должна быть понятна не сама по себе, а как признак эффективности платформы для потоковой обработки данных. Аналогично, производительность криптографических операций должна сопоставляться с типовыми сценариями: скоростью шифрования трафика, интенсивностью хеширования больших файлов или возможным числом асимметричных операций в единицу времени.

Отдельно выделяется требование интерактивного режима криптографических измерений. Пользователь должен иметь возможность выбирать алгоритм, конфигурацию ключа, режим работы и объём данных. Система при этом обязана обеспечивать корректное измерение времени выполнения, учитывать различие между короткими и длинными операциями, а также выводить результат в наглядной форме. Подобная функциональность делает комплекс не только средством наблюдения, но и инструментом исследовательского характера, позволяющим анализировать влияние параметров алгоритма на вычислительные затраты.

Наконец, функциональные требования включают поддержку расширения каталога тестов без переработки всего приложения. Это означает, что добавление нового измерения не должно требовать изменения базовой логики визуализации, механизма запуска или структуры всего интерфейса. Такая постановка задачи непосредственно влияет на архитектурные решения, рассматриваемые далее.

\subsection{Нефункциональные требования}

К ключевым нефункциональным требованиям относятся кросс-платформенность, отзывчивость интерфейса, воспроизводимость результатов, расширяемость и модульность архитектуры. Кросс-платформенность важна по нескольким причинам. Во-первых, она позволяет использовать один и тот же инструмент для анализа различных операционных систем, что существенно повышает прикладную ценность комплекса. Во-вторых, она делает возможным сравнение платформ в сопоставимой программной среде. В-третьих, она снижает стоимость сопровождения, так как вместо нескольких отдельных приложений поддерживается единая кодовая база.

Отзывчивость интерфейса означает, что пользовательский UI не должен блокироваться ни при выполнении серии бенчмарков, ни при периодическом обновлении мониторинговых виджетов. Для инструмента анализа это требование является не косметическим, а принципиальным: если интерфейс теряет отзывчивость во время теста, пользователь лишается возможности наблюдать контекст измерений и взаимодействовать с системой. Следовательно, архитектура должна предусматривать асинхронное выполнение длительных операций и безопасную передачу результатов в слой представления.

Воспроизводимость подразумевает устойчивость результатов к фоновым факторам и осмысленную организацию методики измерений. Она достигается не только наличием повторных прогонов, но и поддержкой разогрева, фильтрации выбросов, оценкой разброса и использованием единых правил расчёта итогового показателя. Без этих механизмов инструмент был бы пригоден лишь для ориентировочной оценки, но не для сравнительного анализа или учебно-исследовательской работы.

Расширяемость и модульность имеют особое значение, поскольку предметная область анализа вычислительных систем постоянно развивается. Появляются новые классы нагрузок, новые криптографические алгоритмы, новые метрики телеметрии и новые сценарии интерпретации результатов. Поэтому набор тестов и отображаемых показателей должен расширяться без изменения базовой архитектуры и без необходимости полностью перерабатывать интерфейс при добавлении нового измерения. Иначе проект быстро утратит сопровождаемость и станет плохо пригоден для дальнейшего развития.

Кроме того, к нефункциональным требованиям можно отнести интерпретируемость, то есть способность системы представлять данные в форме, понятной пользователю. Для прикладного инструмента мало просто собрать показатели; необходимо так организовать их отображение, чтобы пользователь мог связать число с реальным поведением платформы и понять, как использовать этот результат в инженерной или учебной практике.

\section{Проектирование архитектуры и выбор технологического стека}

Проектирование архитектуры программного комплекса начинается с выбора технологического стека, который должен удовлетворять как функциональным, так и нефункциональным требованиям. В данной работе программный комплекс реализован на языке Go, что обусловлено сочетанием практических преимуществ для задач системного анализа. Компиляция в нативный код обеспечивает предсказуемую производительность синтетических тестов и снижает влияние исполняющей среды на характер нагрузки. Это особенно важно для бенчмарков, где нежелательны дополнительные уровни абстракции, способные вносить нестабильность или трудноучитываемые накладные расходы.

Существенным преимуществом Go является развитая стандартная библиотека, включающая основные криптографические примитивы, средства работы со временем, конкурентным исполнением, буферами данных и системным взаимодействием. Для рассматриваемого проекта это означает возможность строить измерительные сценарии на базе хорошо поддерживаемых и переносимых компонентов без чрезмерной зависимости от внешних пакетов. Дополнительным аргументом в пользу Go является модель конкурентности, основанная на горутинах и каналах. Она позволяет естественным образом разделять фоновый сбор телеметрии, выполнение тестов и обновление интерфейса, не прибегая к избыточно сложным механизмам синхронизации на уровне операционной системы.

Для графической части выбрана библиотека Fyne, обеспечивающая единый API для разных платформ и достаточную скорость отрисовки при регулярном обновлении виджетов. В контексте данного проекта выбор GUI-библиотеки определяется не только удобством создания окна и стандартных элементов управления, но и возможностью быстро формировать вкладки, карточки, таблицы, прогресс-индикаторы и интерактивные формы. Поскольку приложение ориентировано на регулярное обновление данных мониторинга, библиотека интерфейса должна выдерживать такой режим без заметного ухудшения пользовательского опыта.

Архитектура построена модульно и отражает границы ответственности:
\begin{itemize}
    \item \texttt{pkg/profiling} --- сбор телеметрии ОС. Модуль получает информацию о модели и характеристиках процессора, использовании памяти, дисковом и сетевом вводе-вывода, средней нагрузке (load average), состоянии дискового пространства, а также формирует список наиболее активных процессов. Для гетерогенных процессоров Apple Silicon дополнительно реализована попытка определения количества P- и E-ядер через вызовы \texttt{sysctl}.
    \item \texttt{pkg/benchmark} --- библиотека тестов и ядро запуска. Все тесты приводятся к единому интерфейсу \texttt{BenchmarkFunc}, что позволяет расширять набор измерений без изменения внешней логики. Запуск осуществляется через \texttt{TestRunner}, который выполняет серию итераций, выполняет «разогрев» для стабилизации состояния, фильтрует выбросы методом IQR и вычисляет робастные показатели (медиана, MAD, коэффициент вариации). Данная схема уменьшает влияние фоновых задач ОС и делает вывод более воспроизводимым.
    \item \texttt{pkg/gui} --- представление данных и управление сценариями пользователя. Точка входа приложения формирует вкладки: «Мониторинг», «Процессор», «Память», «Математика», «Криптография». Панель мониторинга отображает телеметрию с периодом обновления 1 секунду и визуализирует нагрузку по каждому логическому ядру. Панели бенчмарков запускают тесты по выбранной категории и показывают интерпретацию результата в прикладных единицах.
\end{itemize}

Такое разбиение не является случайным; оно отражает принцип разделения ответственности между слоями приложения. Подсистема телеметрии отвечает за получение и предварительную нормализацию данных. Подсистема бенчмарков концентрирует измерительную логику и статистическую обработку. Подсистема GUI отвечает за сценарии взаимодействия с пользователем и визуализацию. Благодаря этому каждый слой остаётся относительно независимым: изменение способов отображения не требует вмешательства в алгоритмы тестирования, а расширение набора метрик не должно приводить к переработке механизмов статистического расчёта.

Подобная архитектура позволяет также изолировать измерительную нагрузку от UI: вычисления выполняются в фоновых горутинах, а интерфейс получает только готовые значения или сообщения о прогрессе. Это поддерживает отзывчивость даже во время серии тестов и уменьшает риск того, что само средство отображения станет значимым источником искажения результатов. Одновременно модульность задаёт понятный путь расширения: новые тесты добавляются как функции в \texttt{pkg/benchmark}, а в UI подключаются через описание задачи и интерпретацию результата без изменения архитектурных слоёв.

\section{Потоки данных и модель обновления интерфейса}

Проектирование комплекса требует явного разделения данных, которые обновляются периодически, и данных, вычисляемых по запросу пользователя. Мониторинговые метрики формируют непрерывный поток обновлений с фиксированным шагом по времени. Они отражают текущее состояние системы и должны отображаться в виде текущих значений, индикаторов и сводных характеристик. Бенчмарки, напротив, запускаются событийно и представляют собой ограниченную по времени серию вычислений, результат которой должен быть представлен пользователю как завершённый измерительный акт вместе со статистикой качества и пояснением смысла показателя.

Такое различие в природе данных непосредственно влияет на архитектуру потоков. Если обращаться с мониторинговыми метриками и результатами тестов одинаково, система либо станет избыточно сложной, либо начнёт смешивать краткоживущие и непрерывные данные в одном интерфейсном контуре. Поэтому в приложении принята модель, при которой сбор телеметрии и выполнение бенчмарков существуют как разные процессы взаимодействия с системой, но сходятся на уровне пользовательского представления.

В практической реализации это означает, что мониторинг выполняется в фоне с предсказуемой периодичностью, а результаты тестов вычисляются в отдельных фоновых задачах. Интерфейс не занимается «сырым» измерением и не должен принимать участие в тяжёлых вычислениях. Его функция заключается в том, чтобы получать готовые агрегированные значения, отображать прогресс, предоставлять пользователю средства запуска сценариев и сохранять визуальную целостность данных.

Подобная модель снижает взаимное влияние мониторинга и измерений. Во-первых, мониторинг остаётся достаточно редким и стабильным, чтобы не искажать результаты тестов чрезмерной собственной активностью. Во-вторых, длительные бенчмарки не блокируют отрисовку интерфейса и не лишают пользователя возможности наблюдать состояние системы во время прогона. В-третьих, сохраняется важный контекст: в момент измерения пользователь может сопоставить итоговый результат с фоновой нагрузкой, активными процессами и поведением системы в реальном времени.

С архитектурной точки зрения можно говорить о двух типах потоков данных: телеметрических и событийно-измерительных. Телеметрические потоки ориентированы на регулярность, малую задержку отображения и устойчивую нормализацию показателей. Событийно-измерительные потоки ориентированы на корректность расчёта, повторяемость методики и на атомарное представление завершённого результата. Разведение этих потоков по логическим уровням позволяет сделать систему одновременно информативной, отзывчивой и устойчивой к росту функциональности.

\section{Унификация измерений и расширяемость каталога тестов}

Для того чтобы пользовательский интерфейс мог одинаково работать с результатами разнородных тестов, все измерения приводятся к единому контракту: тест возвращает числовой показатель и единицу измерения. На первый взгляд такое решение выглядит простым, однако именно оно задаёт важнейшее архитектурное свойство всей системы --- отделение вычислительной части от способа представления результата. Интерфейс не должен знать внутреннюю специфику каждого алгоритма, объём выделяемой памяти, число внутренних итераций или конкретную реализацию теста. Ему достаточно получить числовое значение, единицу измерения и при необходимости связанную функцию интерпретации.

Унификация результатов даёт сразу несколько преимуществ. Во-первых, она снижает связанность между модулем тестирования и GUI, что положительно влияет на сопровождаемость кода. Во-вторых, она упрощает добавление новых тестов: разработчику достаточно реализовать функцию измерения, встроить её в каталог задач и определить способ интерпретации результата. В-третьих, она делает возможным единый механизм статистической обработки, так как все измерения проходят через общую инфраструктуру запуска и агрегации.

Расширяемость каталога тестов особенно важна для предметной области системного анализа, поскольку перечень интересующих пользователя измерений не является фиксированным. Сегодня акцент может быть сделан на CPU, память и криптографию, а в дальнейшем могут потребоваться дополнительные сценарии: файловые профили, сетевые испытания, анализ латентности конкретных структур данных, тесты устойчивости многопоточного выполнения или оценки специализированных ускорителей. Если архитектура не допускает такого роста без значительной переработки, инструмент быстро теряет практическую ценность.

В рассматриваемом проекте добавление нового теста сводится к реализации функции, регистрации задачи в каталоге и подключению описания, необходимого для пользователя. При этом не требуется изменение механизма запуска, дублирование логики статистической обработки или пересборка всей модели представления. Такой подход соответствует принципу архитектурной открытости для расширения и закрытости для модификации, что особенно важно для исследовательских и учебных программных комплексов, развивающихся по мере накопления новых требований.

Следовательно, унификация измерений и продуманная расширяемость каталога тестов выступают не частной технической деталью, а одним из центральных архитектурных решений проекта. Именно они позволяют сохранить логическую целостность системы при росте числа сценариев анализа и обеспечивают возможность дальнейшего развития комплекса без утраты его сопровождаемости и интерпретируемости.
